\documentclass[11pt]{article}
\usepackage[margin=25mm]{geometry}
\usepackage{fontspec}
\setmainfont{Latin Modern Roman}
\usepackage{amsmath,amssymb,mathtools}
\usepackage{unicode-math}
\setmathfont{Latin Modern Math}
\usepackage{xurl,hyperref}
\hypersetup{colorlinks=true,urlcolor=blue}
\setcounter{secnumdepth}{0}
\setlength{\parindent}{0pt}
\setlength{\parskip}{5pt}
\setlength{\emergencystretch}{3em}
\newcommand{\sref}[2]{\hyperlink{source:#1:#2}{[S#2]}}
\newcommand{\cataloguescope}[1]{\paragraph{Recorded target.}#1}
\begin{document}
\section{Research Problems in Function Theory — Problem 6.65}\label{statement}

\subsection{Definitions and mathematical statement}

\paragraph{Shared notation.}$\mathbb N=\{1,2,\ldots\}$, and $\mathbb Z,\mathbb Q,\mathbb R,\mathbb C$ denote the integers, rationals, reals and complex numbers. Any other conventions are specified locally.


Write $\mathbb D=\{z\in\mathbb C:|z|<1\}$ and $\mathbb T=\partial\mathbb D$. The class $S$ consists of holomorphic injective functions $f:\mathbb D\to\mathbb C$ normalized by $f(0)=0$ and $f'(0)=1$, so $f(z)=z+\sum_{n=2}^{\infty}a_nz^n$. For $M>1$, let $S^*(M)$ consist of the functions $f\in S$ with $|f(z)|\le M$ throughout $\mathbb D$ and with image starlike about $0$, meaning $tw\in f(\mathbb D)$ whenever $w\in f(\mathbb D)$ and $0\le t\le1$.
\paragraph{Statement recorded in the supplied source.}
For every real $M$ with $e<M<5$, determine the exact value of $\sup\{|a_3|:f\in S^*(M)\}$. \sref{statement}{2}
\subsection{Short English statement}
Find the sharp third-coefficient bound for bounded starlike functions in the remaining parameter range between e and five.
\subsection{Sources}
\begin{description}
\item[\hypertarget{source:statement:1}{S1}] ULAM.AI and the UnsolvedMath Contributors, \emph{UnsolvedMath: A Curated Collection of Open Mathematics Problems}, catalogue entry. CC BY 4.0. \url{https://huggingface.co/datasets/ulamai/UnsolvedMath}.
\item[\hypertarget{source:statement:2}{S2}] W. K. Hayman and E. F. Lingham, \emph{Research Problems in Function Theory} (2018), Chapter 6, Problem 6.65 and its complete update; Chapter 6 initial notation. \url{https://arxiv.org/abs/1809.07200}.
\end{description}
\label{end:statement}
\end{document}
